\section{Khảo sát các nền tảng Feature Flag hiện có}

Phần này khảo sát hai nền tảng tiêu biểu là LaunchDarkly và Unleash nhằm rút ra các nguyên tắc thiết kế cần kế thừa cho mini project. Nội dung chỉ tập trung vào điểm khác biệt và bài học thiết kế, tránh lặp lại các khái niệm Feature Flag đã trình bày ở phần tổng quan.

\subsection{Nền tảng LaunchDarkly}

LaunchDarkly là nền tảng Feature Flag thương mại dạng dịch vụ được quản lý. Nền tảng tổ chức tài nguyên theo project, environment và flag; cung cấp SDK đa nền tảng; hỗ trợ targeting rule có thứ tự, percentage rollout, nhiều variation, audit log và quản trị quyền truy cập \cite{launchdarkly_feature_flags,launchdarkly_targeting_rules}.

Điểm đáng chú ý đối với mini project là mô hình ordered rule. Người quản trị có thể đặt rule cụ thể trước rule tổng quát, ví dụ user allowlist trước percentage rollout. Cách biểu diễn này giúp kết quả evaluation dễ giải thích vì engine dừng tại rule đầu tiên khớp.

Hạn chế chính khi tham chiếu LaunchDarkly cho mini project là phạm vi tính năng rất rộng. Các khả năng như segment phức tạp, nhiều variation, workflow phê duyệt và phân tích thử nghiệm chuyên sâu được xem là hướng mở rộng, không phải yêu cầu triển khai ban đầu.

\subsection{Nền tảng Unleash}

Unleash là nền tảng Feature Flag mã nguồn mở, có thể self-hosted hoặc dùng dịch vụ được quản lý. Nền tảng sử dụng activation strategy để quyết định flag có được bật hay không trong từng environment; SDK thường tải cấu hình về và đánh giá tại ứng dụng \cite{unleash_feature_flags,unleash_feature_flag_best_practices}.

Điểm đáng chú ý đối với mini project là khả năng tự triển khai và tư duy tách control plane khỏi runtime evaluation. Unleash cũng nhấn mạnh lifecycle của flag, metric evaluation và việc tránh để hệ thống Feature Flag trở thành phụ thuộc đồng bộ bắt buộc của mọi request nghiệp vụ.

Hạn chế của mô hình self-hosted là tổ chức phải tự chịu trách nhiệm về database, backup, monitoring, nâng cấp và mở rộng. Vì mini project cần dễ chạy và dễ trình diễn, hệ thống chỉ mô phỏng phần lõi bằng một backend duy nhất, PostgreSQL và cache đơn giản.

\subsection{So sánh và bài học rút ra}

Bảng~\ref{tab:tong-hop-so-sanh-launchdarkly-unleash} tóm tắt các điểm so sánh có ảnh hưởng trực tiếp đến thiết kế mini project.

\begin{table}[H]
    \centering
    \caption{Tổng hợp so sánh LaunchDarkly và Unleash}
    \label{tab:tong-hop-so-sanh-launchdarkly-unleash}
    \begin{tabular}{p{0.22\textwidth} p{0.30\textwidth} p{0.30\textwidth}}
        \toprule
        \textbf{Tiêu chí} & \textbf{LaunchDarkly} & \textbf{Unleash} \\
        \midrule
        Mô hình triển khai & Dịch vụ được quản lý & Managed hoặc self-hosted \\
        Đơn vị logic & Ordered targeting rule & Activation strategy và constraint \\
        Evaluation & Chủ yếu qua SDK & Chủ yếu qua SDK hoặc proxy/edge \\
        Rollout & Percentage rollout ổn định & Gradual rollout ổn định \\
        Audit/RBAC & Tích hợp phong phú & Phụ thuộc phiên bản/mô hình triển khai \\
        Bài học cho mini project & Precedence rõ, reason dễ giải thích & Tách plane, self-hosted workflow, lifecycle \\
        \bottomrule
    \end{tabular}
\end{table}

Từ khảo sát, mini project kế thừa các nguyên tắc sau:

\begin{itemize}
    \item tổ chức tài nguyên theo project, environment và flag;
    \item tách control plane quản trị khỏi data plane evaluation;
    \item dùng evaluation context thay vì hard-code điều kiện trong ứng dụng;
    \item áp dụng stable rollout bằng khóa không chứa PII;
    \item trả reason code để giải thích kết quả evaluation;
    \item ghi audit log cho mọi mutation quan trọng;
    \item cung cấp SDK và fallback an toàn cho ứng dụng tích hợp.
\end{itemize}

\subsection{Định hướng của hệ thống đề xuất}

Hệ thống đề xuất không cố gắng sao chép đầy đủ LaunchDarkly hoặc Unleash. Phạm vi được chọn là một lát cắt đủ nhỏ để triển khai trong mini project nhưng vẫn thể hiện được bản chất của nền tảng Feature Flag: quản lý cấu hình, đánh giá runtime, thay đổi hành vi ứng dụng, auditability và trình diễn release an toàn.

Các chức năng được kế thừa trực tiếp gồm project/environment, flag config theo environment, ordered rule, user/role targeting, percentage rollout, kill switch, audit log, statistics, SDK và Docker Compose. Các chức năng được đơn giản hóa gồm multi-variant testing, segment phức tạp, approval workflow, realtime streaming, multi-region deployment và phân quyền enterprise.

Điểm khác biệt của MiniProject là ưu tiên tính minh bạch và khả năng giải thích hơn độ rộng tính năng. Vì vậy, Rule Evaluation Engine dùng precedence cố định, response có reason code rõ ràng, audit log append-only và Dashboard tập trung vào các kịch bản demo quan trọng nhất.
